\documentclass[10pt,a4paper]{amsart}
\usepackage[margin=19mm]{geometry}
\usepackage{amsmath,amssymb,mathtools}
\usepackage[scr=rsfs,cal=euler]{mathalfa}
\usepackage{xfrac}
\usepackage[unicode,hidelinks,bookmarks=false]{hyperref}
\newcommand{\ie}{i.e.\ }
\newcommand{\cf}{cf.\ }
\newcommand{\resp}{respectively\ }
\newcommand{\Subsection}{\subsection}
\providecommand{\nobreakdash}{\nobreak\hbox{-}}
\setlength{\emergencystretch}{2em}
\multlinegap0pt
\usepackage{xcolor}
\usepackage{amssymb}
\usepackage[shortlabels]{enumitem}
\usepackage{float}
\usepackage{mathtools}
\newtheorem{lem}{Lemma}
\newcommand{\N}{\mathbb{N}}
\title{Graphical small cancellation and hyperfiniteness of boundary actions}
\author{Chris Karpinski, Damian Osajda, Koichi Oyakawa}
\date{Source: arXiv:2509.05548v2 (CC BY 4.0)}
\begin{document}
\maketitle

\noindent\emph{Source and licence.} Graphical small cancellation and hyperfiniteness of boundary actions. Version arXiv:2509.05548v2 (CC BY 4.0), distributed by arXiv under the Creative Commons Attribution 4.0 International licence, http://creativecommons.org/licenses/by/4.0/, as recorded in the arXiv licence metadata for this exact version and shown on the abstract page https://arxiv.org/abs/2509.05548v2. Full licence notice: \texttt{licenses/arxiv-2509.05548-CC-BY-4.0.txt}.\par\smallskip

\noindent\textit{Editorial scope.} Counted unit: the article's lem closed (source0.tex lines 778--781). The article's complete original proof of it is delivered with it (source0.tex lines 783--800, one \texttt{proof} environment of 18 lines). No mathematics is written, repaired or re-derived: the statement and the proof are the article's own lines, in the article's own notation, and no retained line is substituted or re-flowed.  Retained uncounted as the source-local context the proof needs: lines 502--517; lines 534--538.  Nothing was omitted from any retained span, so the package carries no omission record.  No retained line contains a citation, so the wrapper carries no bibliography and no bibliography entry.  No retained line contains a figure environment, an image-inclusion command or drawing code, so no image is delivered and the package declares no asset dependency.  Editorial formatting only, in addition: an AMS \texttt{amsart} preamble that restates the article's own declarations and macros used by the retained lines, namely the environments \texttt{lem} on separate counters, together with the standard packages those lines need.  The retained environments print in this document as Lemma 1 in order of appearance, whereas the article prints the same items with the numbering of its own full document.  The complete native source of the article as distributed in the arXiv e-print for this exact version is bundled unchanged as \texttt{sources/184578/source0.tex}.
\begin{lem}\label{lem:geod ray converging in Y}
    Let $p=(p(0),p(1),p(2),\ldots)$ be a geodesic ray in $X$, where $p(i)$'s are vertices in $X$ composing $p$. The following conditions (1)-(3) are equivalent.
    \begin{itemize}
        \item[(1)]
        There exists $\xi \in \partial Y$ such that the sequence $(p(n))_{n \in \N}$ converges to $\xi$.
        \item[(2)]
        $\sup_{n \in \N}d_Y(p(0),p(n))= \infty$.
        \item[(3)] 
        There exist $\xi \in \partial Y$ and a sequence of vertices $(a_n)_{n \in \N}$ on $p$ such that (a) and (b) below are satisfied.
    
    \begin{itemize}
        \item[(a)] The subpath of $p$ defined by $r_i = p_{[a_{i-1}, a_i]}$ is either an edge in $X$ not appearing in any relator $\Theta$ or a subpath of $p$ contained in a relator (i.e.\ $r_i$ projects to an edge in $Y$).
        \item[(b)] $a_n \to \xi$ in $\partial Y$ and $(a_n)_n$ defines a geodesic ray in $Y$, i.e.\ $d_Y(a_i, a_j) = \vert i - j \vert$ for all $i,j \in \N$. %the sequence $(\hat{r_i})_{i \in \N}$ of edges in $Y$ defines a geodesic ray in $Y$. 
    \end{itemize}
    \end{itemize}
\end{lem}

\begin{lem}
    \label{lemma: geodesic representatives for boundary}
    For any $\xi \in \partial Y$, we have $G(\xi) \neq \emptyset$.
    
\end{lem}

\begin{lem}
    \label{closed}
    The set $A = \{(\xi, p) \in \partial Y \times S^{\N}: p \in G(\xi)\}$ is closed in $\partial Y \times S^{\N}$.
\end{lem}

\begin{proof}

Let $(\xi_n, p_n)_n \subset A$ converge to $(\xi, \gamma) $ in $\partial Y \times S^{\N}$. 

    We show that there exists a subsequence $(p_{n_k})_k$ of $(p_n)_n$ and a geodesic ray $q \in G(\xi)$ such that $p_{n_k} \to q$.

    Let $p \in G(\xi)$ be arbitrary. For each $n \in \N$, let $(a_m^n)_m \subset p_n$ be a sequence of vertices on $p_n$ as in Lemma \ref{lem:geod ray converging in Y} (3). Similarly, for each $n$, let $(b_n)_n \subset p$ be a sequence of vertices on $p$ and $(r_i)_i$ a sequence of subpaths as in Lemma \ref{lem:geod ray converging in Y} (3). Denote $\Theta_i$ a fixed relator containing \textcolor{black}{$r_i$ or the subpath $r_i$ itself} if $r_i$ is not contained in a relator. 
    
    Fix $n \in \N$. Since $\xi_i \to \xi$, there exists $i \in \N$ such that for all $j > i$, we have $(a_i^j, b_i)_1 > n + \delta + 2$. 

    Letting $\mathcal{A}_n := \{v \in p_j : j > i \text{ and } d_Y(1,v) = n\}$ and arguing as in the proof of Lemma \ref{lemma: geodesic representatives for boundary}, we have that $\mathcal{A}_n \subseteq \Theta_{n-1} \cup \Theta_n \cup \Theta_{n+1}$, hence $\mathcal{A}_n$ is finite. A compactness argument as in the proof of Lemma \ref{lemma: geodesic representatives for boundary} yields that there exists a subsequence of $(p_i)_i$ which converges in $S^{\N}$ to a geodesic ray $q$ in $X$ containing a sequence of vertices $(v_n)_n$ such that for each $n \in \N$, $d_Y(p,v_n) \leq 1$ and $d_Y(1,v_n) = n$ for all $n \in \N$. Therefore, $q \in G(\xi)$. 
    %d_Y(v_m,v_n) = \vert n - m \vert$ for all $n, m \in \N$. 

    Since $p_n \to \gamma$ and a subsequence of $(p_n)_n$ converges to $q$, we have $\gamma = q$. Hence, $\gamma \in G(\xi)$. 

    We conclude that $A$ is closed. 
    
\end{proof}

\end{document}
